\documentclass[10pt]{article}
\usepackage[a4paper,margin=0.55in]{geometry}
\usepackage{amsmath,amssymb,amsfonts,mathtools}
\usepackage[dvipsnames]{xcolor}
\usepackage{array,booktabs,float}
\usepackage[colorlinks=true,linkcolor=blue]{hyperref}
\allowdisplaybreaks
\numberwithin{equation}{section}
\setlength{\jot}{4pt}
\setlength{\emergencystretch}{3em}
\newcommand{\Dh}{\hat\Delta}
\newcommand{\OO}{\mathcal O}
\newcommand{\pre}{\frac{g^{2}}{4\pi^{2}}\frac{N_{c}}{(2\pi)^{3}}\frac{g^{2}}{2\pi^{2}}\frac{1}{k^{+}}}
\newcommand{\preLO}{\frac{1}{(2\pi)^3}\frac{g^2}{2\pi^2k^+}}
\newcommand{\XX}{\frac{(x-y)\cdot(x'-y')}{(x-y)^{2}(x'-y')^{2}}}
\newcommand{\ph}{e^{-ik\cdot(y'-y)\frac{1}{\xi}}}
\newcommand{\phone}{e^{-ik\cdot(y'-y)}}

\title{Group II, Row IV in your form: fragmentation function and DGLAP}
\date{}

\begin{document}
\maketitle
\noindent Throughout:
\begin{itemize}
\item the phase $e^{-ik\cdot(y'-y)/\xi}$, the measure $\frac{d\xi}{\xi^2}$ and the $\frac{1}{k^+}$ are kept explicit;
\item $d_\perp=2-\epsilon$, and the collinear pole is infrared ($\epsilon<0$);
\item $\overline{\rm MS}$ subtracts the pole $-\frac{2}{\epsilon}$ (in my earlier notes $-\frac2\epsilon\equiv\frac{1}{\bar\epsilon}$);
\item $X\equiv\frac{x-y}{(x-y)^2}$, $X'\equiv\frac{x'-y'}{(x'-y')^2}$, $\Dh\equiv\frac{y'-y}{|y'-y|}$, and the colour factor is
\begin{equation*}
\OO\equiv\big[U^{e'c}(y')-U^{e'c}(x')\big]\big[U^{ec}(y)-U^{ce}(x)\big]\rho^{e'}(x')\rho^{e}(x).
\end{equation*}
\end{itemize}

%==========================================================================================
\section{Your form}\label{sec:form}
\begin{align}
\frac{d^{3}N_{\rm IV}}{d^{3}k}={}&\pre\int_{x,y,x',y'}\XX\,W\,\Bigg[\int_{\frac{k^{+}}{\vee}}^{1}\frac{d\xi}{\xi^{2}}\,\ph\Big[\xi(1-\xi)+\frac{1-\xi}{\xi}+\frac{\xi}{(1-\xi)_{+}}\Big]\Bigg]\OO\notag\\
&+\pre\Big[\log\frac{\vee}{\Lambda}+\log\Big(1-\frac{k^{+}}{\vee}\Big)\Big]\int_{x,y,x',y'}\phone\,\XX\,W\,\OO\notag\\
&+\pre\int_{\frac{k^{+}}{\vee}}^{\frac{k^{+}}{k^{+}+\Lambda}}d\xi\,\frac{1-\xi}{\xi}\int_{x,y,x',y'}\ph\notag\\
&\qquad\times\Big[\XX-d_\perp\frac{(x-y)\cdot(y'-y)\,(x'-y')\cdot(y'-y)}{(x-y)^{2}(x'-y')^{2}(y'-y)^{2}}\Big]\OO ,
\label{eq:form}
\end{align}
with the weight
\begin{equation}
W\equiv-\frac{d_{\perp}}{\epsilon}-d_{\perp}\gamma_{E}-\frac{d_{\perp}}{2}\log\frac{(y'-y)^{2}\mu^{2}_{\overline{MS}}}{4}-1 ,
\end{equation}
and your subtraction at $\xi=1$:
\begin{equation}
\int_{\frac{k^{+}}{\vee}}^{1}\frac{d\xi}{\xi^{2}}\,\ph\Big[\dots+\frac{\xi}{(1-\xi)_{+}}\Big]\equiv\int_{\frac{k^{+}}{\vee}}^{1}d\xi\bigg\{\frac{\ph-\phone}{\xi(1-\xi)}+\Big[\xi(1-\xi)+\frac{1-\xi}{\xi}\Big]\frac{1}{\xi^2}\,\ph\bigg\}.
\label{eq:usub}
\end{equation}
\begin{itemize}
\item \textbf{The $\delta(1-\xi)$ coefficient is exact:}
\begin{equation*}
\int_{k^+/\vee}^{k^+/(k^++\Lambda)}\frac{d\xi}{\xi(1-\xi)}=\log\frac{\vee}{\Lambda}+\log\Big(1-\frac{k^+}{\vee}\Big)\qquad(\texttt{sympy}: $0$).
\end{equation*}
\item \textbf{The upper limit 1} instead of $\frac{k^+}{k^++\Lambda}$ in the first line adds the strip $\frac{k^+}{k^++\Lambda}<\xi<1$, i.e.\ $\hat R=-i\frac{\Lambda}{k^+}k\cdot(y'-y)\,e^{-ik\cdot(y'-y)}+O(\frac{\Lambda^2}{k^{+2}})$ inside the braces (times $W$ and the colour structure).
\begin{itemize}
\item With upper limit $\frac{k^+}{k^++\Lambda}$, Eq.~\eqref{eq:form} equals Row IV to rel.\ $2\times10^{-14}$.
\item With upper limit 1, at $\frac{\Lambda}{k^+}=10^{-4}$, the difference is $1.3\times10^{-4}$.
\end{itemize}
Below, $\Lambda\to0$ is taken only in such power terms. $\log\frac{\vee}{\Lambda}$ is kept.
\end{itemize}

%==========================================================================================
\section{Step 1: pole and finite part of the weight}\label{sec:W}
With $d_\perp=2-\epsilon$, $-\frac{d_\perp}{\epsilon}=-\frac{2}{\epsilon}+1$. The $+1$ cancels the $-1$ of $W$:
\begin{equation}
W=-\frac{2}{\epsilon}-2\gamma_E-\log\frac{(y'-y)^2\mu^2_{\overline{MS}}}{4}+O(\epsilon)
\label{eq:W}
\end{equation}
(\texttt{sympy}: difference $0$). For $d_\perp=2$ instead, $W=-\frac2\epsilon-2\gamma_E-\log\frac{(y'-y)^2\mu^2}{4}-1$. The pole is the same, and the finite part changes by a scheme term.
\begin{itemize}
\item The pole of Row IV is the first two lines of Eq.~\eqref{eq:form} with $W\to-\frac2\epsilon$.
\item The third line has no pole.
\item With $\frac{g^2N_c}{4\pi^2}=\frac{\alpha_s}{2\pi}\cdot2N_c$ (\texttt{sympy}: $0$), the pole is
\end{itemize}
\begin{align}
\frac{d^{3}N_{\rm IV}}{d^{3}k}\bigg|_{\rm pole}={}&\frac{\alpha_s}{2\pi}\Big(-\frac{2}{\epsilon}\Big)\preLO\int_{x,y,x',y'}\XX\,\OO\notag\\
&\times2N_c\bigg\{\int_{\frac{k^{+}}{\vee}}^{1}d\xi\Big[\frac{\ph-\phone}{\xi(1-\xi)}+\Big(\xi(1-\xi)+\frac{1-\xi}{\xi}\Big)\frac{\ph}{\xi^2}\Big]\notag\\
&\qquad+\Big[\log\frac{\vee}{\Lambda}+\log\Big(1-\frac{k^{+}}{\vee}\Big)\Big]\phone\bigg\}.
\label{eq:poleu}
\end{align}

%==========================================================================================
\section{Step 2: from your subtraction to the standard plus distribution}\label{sec:conv}
\paragraph{Why this step is needed.} A fragmentation function $D(\xi,\mu)$ and its DGLAP kernel depend on $\xi$ only.
\begin{itemize}
\item In your split, the subtraction term of Eq.~\eqref{eq:usub} is $\phone\int_{k^+/\vee}^1\frac{d\xi}{\xi(1-\xi)}$. Its $\frac1\xi$ part gives $\int_{k^+/\vee}^1\frac{d\xi}{\xi}=\log\frac{\vee}{k^+}$, which depends on the lower limit $\frac{k^+}{\vee}$.
\item Correspondingly, your $\delta(1-\xi)$ coefficient contains $\log\frac{\vee}{k^+}+\log(1-\frac{k^+}{\vee})=\log\frac{\vee-k^+}{k^+}$, which depends on $k^+/\vee$.
\item In the convolution with $D$ the gluon momentum is $\frac{k^+}{\zeta}$, so the lower limit becomes $\frac{k^+}{\zeta\vee}$, and this piece changes with $\zeta$. It therefore cannot be part of a kernel $P(\xi)$.
\end{itemize}
The standard plus distribution on $[0,1]$ subtracts only $\frac{\phone}{1-\xi}$:
\begin{align}
\int_{\frac{k^+}{\vee}}^1\frac{d\xi}{\xi^2}P^{(r)}_{gg,+}(\xi)\,\ph\equiv{}&2N_c\bigg\{\int_{\frac{k^{+}}{\vee}}^{1}d\xi\Big[\frac{\frac1\xi\ph-\phone}{1-\xi}+\Big(\xi(1-\xi)+\frac{1-\xi}{\xi}\Big)\frac{\ph}{\xi^2}\Big]\notag\\
&+\log\Big(1-\frac{k^+}{\vee}\Big)\phone\bigg\},
\label{eq:std}
\end{align}
with $P^{(r)}_{gg,+}(\xi)=2N_c\big[\frac{\xi}{(1-\xi)_+}+\frac{1-\xi}{\xi}+\xi(1-\xi)\big]$ and $\frac{\xi}{(1-\xi)_+}\equiv\xi\big[\frac{1}{1-\xi}\big]_+$. The $\log(1-\frac{k^+}{\vee})$ is produced by the plus distribution itself, because the integration starts at $\frac{k^+}{\vee}$.

\paragraph{The conversion (exact).} The two integrands differ by
\begin{equation*}
\frac{\ph-\phone}{\xi(1-\xi)}-\frac{\frac1\xi\ph-\phone}{1-\xi}=-\phone\Big[\frac{1}{\xi(1-\xi)}-\frac{1}{1-\xi}\Big]=-\frac{\phone}{\xi},
\end{equation*}
and $\int_{k^+/\vee}^1\frac{d\xi}{\xi}=\log\frac{\vee}{k^+}$. Hence
\begin{align}
&\int_{\frac{k^{+}}{\vee}}^{1}d\xi\Big[\frac{\ph-\phone}{\xi(1-\xi)}+\Big(\xi(1-\xi)+\frac{1-\xi}{\xi}\Big)\frac{\ph}{\xi^2}\Big]+\Big[\log\frac{\vee}{\Lambda}+\log\Big(1-\frac{k^{+}}{\vee}\Big)\Big]\phone\notag\\
&\qquad=\frac{1}{2N_c}\int_{\frac{k^+}{\vee}}^1\frac{d\xi}{\xi^2}P^{(r)}_{gg,+}(\xi)\,\ph+\Big[\log\frac{\vee}{\Lambda}-\log\frac{\vee}{k^+}\Big]\phone
\label{eq:conv}
\end{align}
(numerically, 20 random cases, rel.\ $1\times10^{-15}$). The shift between the two splits is $\phone\big[\log\frac{\vee}{k^+}+\log(1-\frac{k^+}{\vee})\big]$ (\texttt{sympy}: twice it is $2\log\frac{\vee}{k^+}+2\log(1-\frac{k^+}{\vee})$). It moves from your $\delta(1-\xi)$ term into the plus integral.

\paragraph{The pole of Row IV in standard form.}
\begin{equation}
\boxed{\begin{aligned}\frac{d^{3}N_{\rm IV}}{d^{3}k}\bigg|_{\rm pole}={}&\frac{\alpha_s}{2\pi}\Big(-\frac{2}{\epsilon}\Big)\preLO\int_{x,y,x',y'}\XX\,\OO\\
&\times\bigg\{\int_{\frac{k^+}{\vee}}^1\frac{d\xi}{\xi^2}P^{(r)}_{gg,+}(\xi)\,\ph+2N_c\Big[\log\frac{\vee}{\Lambda}-\log\frac{\vee}{k^+}\Big]\phone\bigg\}\end{aligned}}
\label{eq:polestd}
\end{equation}
\begin{itemize}
\item The first term is the real part of $P_{gg}$, convolved with the phase $\ph$ and the measure $\frac{d\xi}{\xi^2}$.
\item The second is the soft logarithm $2N_c\log\frac{k^+}{\Lambda}$ at $\xi=1$ (phase without $\frac1\xi$).
\end{itemize}

%==========================================================================================
\section{Step 3: the \texorpdfstring{$\delta(1-\xi)$}{delta(1-xi)} terms with Group I}\label{sec:delta}
Terms with $\phone$ (no $\frac1\xi$) and the LO structure also come from Group I: Row II and Row V (2b, UV), each with its c.c. There is also the charge renormalization of the LO coupling, $\alpha_0=\alpha_s(\mu)\big[1+\frac{\alpha_s\beta_0}{4\pi}(-\frac2\epsilon)\big]$, with $\beta_0=\frac{11}{3}N_c$. For $\Lambda\to0$ in power terms, the coefficients of
\begin{equation*}
\frac{\alpha_s}{2\pi}\Big(-\frac{2}{\epsilon}\Big)N_c\preLO\int_{x,y,x',y'}\phone\,\XX\,\OO
\end{equation*}
are:
\begin{table}[H]\centering\small
\renewcommand{\arraystretch}{1.3}
\begin{tabular}{@{}lll@{}}\toprule
source & weight & coefficient\\\midrule
Row IV, Eq.~\eqref{eq:polestd} (standard split) & 1 & $2\log\frac{\vee}{\Lambda}-2\log\frac{\vee}{k^+}$\\
\quad (in your split, Eq.~\eqref{eq:poleu}) & & \quad ($2\log\frac{\vee}{\Lambda}+2\log(1-\frac{k^+}{\vee})$, with the shift in the plus integral)\\
Group I, Row II & $\times2$ & $-\frac{11}{3}+2\log\frac{\vee}{\Lambda}-2\log\frac{\vee}{k^+}-\log\big(1-\frac{k^+}{\vee}\big)$\\
Group I, Row V 2b UV & $\times2$ & $\frac{11}{3}-4\log\frac{\vee}{\Lambda}+4\log\frac{\vee}{k^+}$\\
charge renormalization & & $+\frac{11}{6}$\\\midrule
sum & & $\frac{11}{6}-\log\big(1-\frac{k^+}{\vee}\big)$ \quad(\texttt{sympy}: $\frac{11}{6}-\log(1-\xi_0)$)\\\bottomrule
\end{tabular}
\end{table}
\begin{itemize}
\item $\log\frac{\vee}{\Lambda}$: $+2+2-4=0$, and $\log\frac{\vee}{k^+}$: $-2-2+4=0$. The soft logarithms cancel, as they must.
\item The $\frac{11}{3}$ cancel between Row II and Row V 2b.
\item What is left, $\frac{11N_c}{6}=\frac{\beta_0}{2}$, is the $\delta(1-\xi)$ coefficient of $P_{gg}$. The momentum sum rule $\int_0^1d\xi\,\xi P_{gg}=0$ gives the same value (\texttt{sympy}: $\frac{11N_c}{6}$).
\item \textbf{Remainder:} $-N_c\log(1-\frac{k^+}{\vee})$, of order $\frac{k^+}{\vee}$. It depends on $k^+/\vee$, not on $\xi$ alone, so it cannot go into $D$. It must cancel against other rows, or the upper limits $\vee$ of Row II and Row IV must be made consistent; this is still open.
\end{itemize}
So, up to the remainder, the full pole of the gluon spectrum is
\begin{equation}
\frac{d^3N_g}{d^3k}\bigg|_{\rm pole}=\frac{\alpha_s}{2\pi}\Big(-\frac{2}{\epsilon}\Big)\preLO\int_{\frac{k^+}{\vee}}^1\frac{d\xi}{\xi^2}\,P_{gg}(\xi)\int_{x,y,x',y'}\ph\,\XX\,\OO ,
\label{eq:totpole}
\end{equation}
\begin{equation}
P_{gg}(\xi)=2N_c\Big[\frac{\xi}{(1-\xi)_+}+\frac{1-\xi}{\xi}+\xi(1-\xi)\Big]+\frac{11N_c}{6}\,\delta(1-\xi)\qquad(n_f=0).
\label{eq:Pgg}
\end{equation}
The $\delta(1-\xi)$ term sets $\xi=1$ in the phase and in $\frac{1}{\xi^2}$.

%==========================================================================================
\section{Step 4: the fragmentation function}\label{sec:FF}
\paragraph{Factorization, written out.} A final state $h$ with $(k^+,\mathbf k)$ comes from a gluon with $(q^+,\mathbf q)=(\frac{k^+}{\xi},\frac{\mathbf k}{\xi})$:
\begin{equation*}
\frac{d^3N_h}{d^3k}=\int_{\frac{k^+}{\vee}}^1d\xi\,D^{(0)}_{h/g}(\xi)\int dq^+d^2\mathbf q\,\delta(k^+-\xi q^+)\,\delta^{(2)}(\mathbf k-\xi\mathbf q)\,\frac{d^3N_g}{d^3q}(q^+,\mathbf q).
\end{equation*}
At LO, $\frac{d^3N_g}{d^3q}=\frac{1}{(2\pi)^3}\frac{g^2}{2\pi^2q^+}\int e^{-i\mathbf q\cdot(y'-y)}X\cdot X'\,\OO$. Under the substitution:
\begin{itemize}
\item the $\delta$ functions give $\frac{1}{\xi^3}$;
\item $\frac{1}{q^+}=\frac{\xi}{k^+}$;
\item $e^{-i\mathbf q\cdot(y'-y)}=\ph$.
\end{itemize}
So
\begin{equation}
\frac{d^3N^{\rm LO}_h}{d^3k}=\preLO\int_{\frac{k^+}{\vee}}^1\frac{d\xi}{\xi^2}\,D^{(0)}_{h/g}(\xi)\int_{x,y,x',y'}\ph\,\XX\,\OO .
\label{eq:LOh}
\end{equation}
This has the same $\frac{d\xi}{\xi^2}$, phase and $\frac{1}{k^+}$ as Eq.~\eqref{eq:totpole}.

\paragraph{The counterterm.}
\begin{equation}
\boxed{\begin{aligned}&\Gamma_{gg}(z)=\delta(1-z)+\frac{\alpha_s}{2\pi}\,\frac{2}{\epsilon}\,P_{gg}(z),\\
&D^{(0)}_{h/g}(\xi)=\int_\xi^1\frac{dz}{z}\,\Gamma_{gg}(z)\,D_{h/g}\Big(\frac{\xi}{z},\mu\Big)=D_{h/g}(\xi,\mu)+\frac{\alpha_s}{2\pi}\frac{2}{\epsilon}\int_\xi^1\frac{dz}{z}P_{gg}(z)D_{h/g}\Big(\frac{\xi}{z},\mu\Big).\end{aligned}}
\label{eq:Gamma}
\end{equation}
At $O(\alpha_s)$ there are two pole terms.
\begin{itemize}
\item \textbf{A: Eq.~\eqref{eq:Gamma} inserted into Eq.~\eqref{eq:LOh}:}
\begin{equation*}
A=\frac{\alpha_s}{2\pi}\frac{2}{\epsilon}\preLO\int_{\frac{k^+}{\vee}}^1\frac{d\xi}{\xi^2}\Big[\int_\xi^1\frac{dz}{z}P_{gg}(z)D_{h/g}\Big(\frac{\xi}{z},\mu\Big)\Big]\int_{x,y,x',y'}\ph\,\XX\,\OO .
\end{equation*}
\item \textbf{B: the pole \eqref{eq:totpole}, at the gluon momentum $(\frac{k^+}{\zeta},\frac{\mathbf k}{\zeta})$, fragmented with $\int\frac{d\zeta}{\zeta^3}D_{h/g}(\zeta,\mu)$.} In Eq.~\eqref{eq:totpole}:
\begin{itemize}
\item $\frac{1}{k^+}\to\frac{\zeta}{k^+}$;
\item the lower limit becomes $\frac{k^+}{\zeta\vee}$;
\item the integration variable is renamed $\xi\to z$;
\item the phase becomes $e^{-ik\cdot(y'-y)/(\zeta z)}$.
\end{itemize}
Then
\begin{equation*}
B=\frac{\alpha_s}{2\pi}\Big(-\frac{2}{\epsilon}\Big)\preLO\int_{\frac{k^+}{\vee}}^1\frac{d\zeta}{\zeta^2}D_{h/g}(\zeta,\mu)\int_{\frac{k^+}{\zeta\vee}}^1\frac{dz}{z^2}P_{gg}(z)\int_{x,y,x',y'}e^{-ik\cdot(y'-y)\frac{1}{\zeta z}}\,\XX\,\OO .
\end{equation*}
\end{itemize}
In $B$ put $\xi=\zeta z$ at fixed $z$:
\begin{itemize}
\item the phase $e^{-ik\cdot(y'-y)/(\zeta z)}=\ph$;
\item the measure $\frac{d\zeta}{\zeta^2}\frac{dz}{z^2}=\frac{d\xi}{\xi^2}\frac{dz}{z}$;
\item $D_{h/g}(\zeta)=D_{h/g}(\frac{\xi}{z})$;
\item the limits become $\frac{k^+}{\vee}\le\xi\le1$, $\xi\le z\le1$.
\end{itemize}
So $B=-A$. Numerically this was checked with the phase and the plus distribution, including the $\log(1-\frac{k^+}{\zeta\vee})$ and $\log(1-\xi)$ from the cut supports: rel.\ $1\times10^{-12}$. The pole is absorbed into $D_{h/g}$.

%==========================================================================================
\section{Step 5: Row IV after the subtraction, in your form}\label{sec:fin}
Eq.~\eqref{eq:conv} holds for any weight that does not depend on $\xi$. So the subtracted Row IV can be written in your split; only $W\to W+\frac{2}{\epsilon}$ changes:
\begin{align}
\frac{d^{3}N^{\overline{MS}}_{\rm IV}}{d^{3}k}={}&\pre\int_{x,y,x',y'}\XX\Big[-2\gamma_E-\log\frac{(y'-y)^2\mu^2_{\overline{MS}}}{4}\Big]\notag\\
&\qquad\times\int_{\frac{k^{+}}{\vee}}^{1}d\xi\Big[\frac{\ph-\phone}{\xi(1-\xi)}+\Big(\xi(1-\xi)+\frac{1-\xi}{\xi}\Big)\frac{\ph}{\xi^2}\Big]\OO\notag\\
&+\pre\Big[\log\frac{\vee}{\Lambda}+\log\Big(1-\frac{k^{+}}{\vee}\Big)\Big]\notag\\
&\qquad\times\int_{x,y,x',y'}\phone\,\XX\Big[-2\gamma_E-\log\frac{(y'-y)^2\mu^2_{\overline{MS}}}{4}\Big]\OO\notag\\
&+\pre\int_{\frac{k^{+}}{\vee}}^{\frac{k^{+}}{k^{+}+\Lambda}}d\xi\,\frac{1-\xi}{\xi}\int_{x,y,x',y'}\ph\notag\\
&\qquad\times\Big[\XX-2\frac{(x-y)\cdot(y'-y)\,(x'-y')\cdot(y'-y)}{(x-y)^{2}(x'-y')^{2}(y'-y)^{2}}\Big]\OO .
\label{eq:fin}
\end{align}
\begin{itemize}
\item For $d_\perp=2$, each $\big[-2\gamma_E-\log\frac{(y'-y)^2\mu^2}{4}\big]$ becomes $\big[-2\gamma_E-\log\frac{(y'-y)^2\mu^2}{4}-1\big]$, and the last line keeps $d_\perp=2$.
\item The logarithm $-2\gamma_E-\log\frac{(y'-y)^2\mu^2}{4}=-\log\frac{(y'-y)^2\mu^2}{c_0^2}$, with $c_0=2e^{-\gamma_E}$, vanishes at $\mu=c_0/|y'-y|$, i.e.\ $\mu\sim|\mathbf k|$.
\end{itemize}

%==========================================================================================
\section{Step 6: DGLAP}\label{sec:dglap}
\paragraph{From the $\mu$ independence of the cross section.}
\begin{itemize}
\item In Eq.~\eqref{eq:fin}, $\frac{\partial}{\partial\log\mu^2}\big[-2\gamma_E-\log\frac{(y'-y)^2\mu^2}{4}\big]=-1$. So Row IV contributes $-1$ times Eq.~\eqref{eq:poleu} without $(-\frac{2}{\epsilon})$, i.e.\ by Eq.~\eqref{eq:conv} $-1$ times the braces of Eq.~\eqref{eq:polestd}.
\item Group I Rows II and V 2b contribute their $\delta(1-\xi)$ coefficients with $-1$ (their poles also come as $-\frac2\epsilon-2\gamma_E-\log(\dots\mu^2)$).
\item The running of $g^2(\mu)$ in LO gives $-\frac{\alpha_s}{2\pi}\frac{\beta_0}{2}$.
\item By Sec.~\ref{sec:delta} the sum is $-P_{gg}$:
\begin{equation*}
\frac{\partial}{\partial\log\mu^2}\frac{d^3N^{\overline{MS}}_g}{d^3k}=-\frac{\alpha_s}{2\pi}\preLO\int_{\frac{k^+}{\vee}}^1\frac{d\xi}{\xi^2}P_{gg}(\xi)\int_{x,y,x',y'}\ph\,\XX\,\OO .
\end{equation*}
\end{itemize}
Fragment this with $\int\frac{d\zeta}{\zeta^3}D_{h/g}(\zeta,\mu)$ at the gluon momentum $(\frac{k^+}{\zeta},\frac{\mathbf k}{\zeta})$, as for $B$. Add $\frac{\partial D}{\partial\log\mu^2}$ in Eq.~\eqref{eq:LOh}. Then $\frac{\partial}{\partial\log\mu^2}\frac{d^3N_h}{d^3k}=0$ reads
\begin{equation}
\preLO\int_{\frac{k^+}{\vee}}^1\frac{d\xi}{\xi^2}\Big[\frac{\partial D_{h/g}(\xi,\mu)}{\partial\log\mu^2}-\frac{\alpha_s}{2\pi}\int_\xi^1\frac{dz}{z}P_{gg}(z)D_{h/g}\Big(\frac{\xi}{z},\mu\Big)\Big]\int_{x,y,x',y'}\ph\,\XX\,\OO=0 .
\end{equation}
This holds for every $k^+$, every $\mathbf k$, and every projectile and target. Different $\xi$ correspond, through $\ph$, to different parent momenta $\mathbf k/\xi$. So the bracket vanishes for each $\xi$:
\begin{equation}
\boxed{\mu^2\frac{\partial}{\partial\mu^2}D_{h/g}(\xi,\mu)=\frac{\alpha_s(\mu)}{2\pi}\int_\xi^1\frac{dz}{z}\,P_{gg}(z)\,D_{h/g}\Big(\frac{\xi}{z},\mu\Big),\qquad P_{gg}\ \text{of Eq.~\eqref{eq:Pgg}}.}
\label{eq:DGLAP}
\end{equation}

\paragraph{The same from $D^{(0)}$.} In $d_\perp=2-\epsilon$, $\frac{d\alpha_s}{d\log\mu^2}=-\frac{\epsilon}{2}\alpha_s+O(\alpha_s^2)$. So
\begin{equation*}
\frac{\partial}{\partial\log\mu^2}\Big[\frac{\alpha_s}{2\pi}\frac{2}{\epsilon}\Big]=-\frac{\alpha_s}{2\pi}\qquad(\texttt{sympy}: $1$),
\end{equation*}
and $\frac{dD^{(0)}}{d\log\mu^2}=0$ in Eq.~\eqref{eq:Gamma} gives Eq.~\eqref{eq:DGLAP}.

\paragraph{The fragmentation function to first order.}
\begin{equation}
D_{h/g}(\xi,\mu)=D_{h/g}(\xi,\mu_0)+\frac{\alpha_s}{2\pi}\log\frac{\mu^2}{\mu_0^2}\int_\xi^1\frac{dz}{z}P_{gg}(z)\,D_{h/g}\Big(\frac{\xi}{z},\mu_0\Big)+O(\alpha_s^2).
\end{equation}

%==========================================================================================
\section{Checks}
Script \texttt{user\_ff\_checks.py}, results in \texttt{user\_ff\_checks.json}. All numerical checks keep the complex phase.
\begin{center}\small
\begin{tabular}{@{}p{0.64\textwidth}p{0.3\textwidth}@{}}\toprule
check & result\\\midrule
your form, upper limit $\frac{k^+}{k^++\Lambda}$, vs.\ the exact Row IV (9 cases) & rel.\ $2\times10^{-14}$\\
your form, upper limit 1: difference $\to-i\frac{\Lambda}{k^+}k\cdot(y'-y)e^{-ik\cdot(y'-y)}$ & $1.3\times10^{-4}$ at $\Lambda/k^+=10^{-4}$\\
$\int\frac{d\xi}{\xi(1-\xi)}$ between the limits (\texttt{sympy}) & $0$\\
$W$ at $d_\perp=2-\epsilon$, Eq.~\eqref{eq:W} (\texttt{sympy}) & $0$\\
$\frac{g^2N_c}{4\pi^2}-\frac{\alpha_s}{2\pi}2N_c$ (\texttt{sympy}) & $0$\\
your subtraction $\to$ standard plus, Eq.~\eqref{eq:conv} & rel.\ $1\times10^{-15}$\\
$\delta(1-\xi)$ tally (\texttt{sympy}) & $\frac{11}{6}-\log(1-\xi_0)$\\
momentum sum rule (\texttt{sympy}) & $\frac{11N_c}{6}$\\
$B=-A$ with the phase, plus distribution and cut supports & rel.\ $1\times10^{-12}$\\
$\partial_{\log\mu^2}\big[\frac{\alpha_s}{2\pi}\frac{2}{\epsilon}\big]$ over $-\frac{\alpha_s}{2\pi}$ (\texttt{sympy}) & $1$\\\bottomrule
\end{tabular}
\end{center}
\end{document}
